Consider a vortex within the fluid. Its vorticity is carried by the flow
and changes through stretching and viscous diffusion:
\[
 \partial_t\omega+(u\cdot\nabla)\omega
 =(\omega\cdot\nabla)u+\nu\Delta\omega,
 \qquad \omega=\nabla\times u.
\]
Follow a tube of fluid particles along its core, where rotation is
concentrated. If strain stretches the tube at an axial rate \(\sigma>0\),
its length \(\ell\) grows as
\[
 \dot\ell=\sigma\ell>0.
\]
Because the fluid is incompressible, the tube must narrow as it lengthens.
For a locally circular tube enclosing circulation \(\Gamma\), its radius
\(r\) and mean swirl speed \(v\) satisfy
\[
 r\propto\ell^{-1/2},
 \qquad v=\frac{|\Gamma|}{2\pi r}.
\]
At fixed circulation, this narrowing makes the fluid circulate faster.
But viscosity is acting throughout the narrowing. It diffuses momentum
between neighbouring regions, reducing velocity differences and spreading
rotation through the surrounding fluid. Its diffusion distance
\(r_{\mathrm{diff}}\) grows with time, and its contribution to the change
in total squared vorticity is nonpositive:
\[
 r_{\mathrm{diff}}\sim\sqrt{\nu t},
 \qquad
 \langle\omega,\nu\Delta\omega\rangle
 =-\nu\|\nabla\omega\|_2^2\le0.
\]
This spreading can carry rotation beyond the tube's boundary even while
the fluid particles within it move inward. The tube can narrow while
losing enclosed circulation, so the swirl speeds up only if that circulation
does not weaken too quickly.
